\documentclass[10pt,letterpaper]{article}
\usepackage{resumen}
\usetikzlibrary{babel}
\DeclareSIUnit\AA{\mbox{\AA}}
\DeclareSIUnit\erg{erg}
\begin{document}
\titulo{Fisicoquímica de superficies}{Tensión superficial, Young--Laplace, Kelvin, capilaridad y adsorción de Gibbs}

%------------------------------------------------------------------ 1
\section{Interfase y tensión superficial}
La \textbf{interfase} es la zona tridimensional, de \SIrange{10}{100}{\AA}, que separa dos fases. Dentro de una fase las atracciones sobre una molécula se anulan; en la interfase queda una \emph{fuerza neta hacia el interior}, de modo que el sistema tiende a minimizar su superficie. La \textbf{tensión superficial} $\gamma$ es la energía necesaria para aumentar la superficie en una unidad de área: mide las fuerzas internas que se deben vencer.

\begin{minipage}[t]{0.47\linewidth}\vspace{0pt}
\centering
\begin{tikzpicture}[scale=0.82,>=Stealth]
  \fill[accsoft] (0,0) rectangle (5.2,2.6);
  \draw[acc,line width=1.2pt] (0,2.6)--(5.2,2.6);
  \node[muted,font=\scriptsize,anchor=south west] at (0.05,2.65) {vapor};
  \node[muted,font=\scriptsize,anchor=north west] at (0.05,2.55) {líquido};
  % molécula interior
  \fill[acc] (1.3,1.45) circle (0.13);
  \foreach \a in {0,45,...,315}{\draw[->,acc,thick] (1.3,1.45)++(\a:0.2)--++(\a:0.55);}
  % molécula de superficie
  \fill[acc] (3.9,2.45) circle (0.13);
  \foreach \a in {180,0,-45,-135}{\draw[->,acc,thick] (3.9,2.45)++(\a:0.2)--++(\a:0.55);}
  \draw[->,bad,line width=1.6pt] (3.9,2.3)--(3.9,1.0);
  \node[bad,font=\scriptsize,anchor=west] at (4.0,1.35) {fuerza neta};
  \node[font=\scriptsize,align=center] at (1.3,0.2) {interior: fuerzas se anulan};
\end{tikzpicture}\\[-2pt]
{\small Fuerzas sobre una molécula en el interior y en la superficie.}
\end{minipage}\hfill
\begin{minipage}[t]{0.50\linewidth}\vspace{0pt}
\centering
\begin{tikzpicture}
\begin{axis}[resumen,width=\linewidth,height=3.9cm,xlabel={$T$ (K)},ylabel={$\gamma$ (\si{\milli\newton\per\meter})},
  xmin=280,xmax=400,ymin=8,ymax=26,xtick={300,340,380},ytick={10,15,20,25}]
  \addplot[teal,domain=280:400,samples=40]{22.0*((1-x/513.9)/(1-298.15/513.9))^(11/9)};
  \addplot[only marks,mark=*,bad,mark size=1.8pt] coordinates {(298.15,22.0)};
\end{axis}
\end{tikzpicture}\\[-2pt]
{\small $\gamma$ del etanol disminuye al subir $T$.}
\end{minipage}

\subsection*{Formulación termodinámica}
A $P$ y $T$ constantes la superficie es trabajo almacenado como energía de Gibbs. Como $\gamma>0$, el área disminuye de forma espontánea:
\begin{clave}
$\displaystyle dG=\gamma\,d\mathcal{A}\qquad\gamma=\left(\frac{\partial G}{\partial\mathcal{A}}\right)_{T,P}>0\qquad \Delta G=\gamma\,(\mathcal{A}_2-\mathcal{A}_1)$
\end{clave}

\begin{minipage}[t]{0.40\linewidth}\vspace{0pt}
\centering\small
\begin{tabular}{@{}ccl@{}}\toprule
$\Delta G$ & $\Delta\mathcal{A}$ & Proceso\\\midrule
$>0$ & $\mathcal{A}_2>\mathcal{A}_1$ & No espontáneo\\
$=0$ & $\mathcal{A}_2=\mathcal{A}_1$ & Equilibrio\\
$<0$ & $\mathcal{A}_2<\mathcal{A}_1$ & Espontáneo\\\bottomrule
\end{tabular}
\end{minipage}\hfill
\begin{minipage}[t]{0.57\linewidth}\vspace{0pt}
\centering\small
\begin{tabular}{@{}lcl@{}}\toprule
Líquido & $\gamma$ (\si{\dina\per\centi\meter}) & Interacción dominante\\\midrule
Hg & 476 & enlace metálico\\
H$_2$O & 72,8 & enlace de hidrógeno\\
C$_6$H$_6$ & 28,9 & interacción entre electrones $\pi$\\
Heptano & 19,7 & fuerzas de dispersión\\\bottomrule
\end{tabular}
\end{minipage}

\smallskip
\textbf{Unidades:} \SI{1}{\dina\per\centi\meter} = \SI{1}{\milli\newton\per\meter} y \SI{1}{\erg\per\centi\meter\squared} = \SI{1}{\milli\joule\per\meter\squared}.

\subsection*{Dependencia con la temperatura}
$\gamma$ disminuye al subir $T$ y se anula en la temperatura crítica $T_c$. Katayama--Guggenheim, con $\gamma_0$ constante del líquido y $n=11/9\approx1{,}22\approx1$ para orgánicos:
\begin{clave}
$\displaystyle \gamma=\gamma_0\left(1-\frac{T}{T_c}\right)^{n}\qquad\Rightarrow\qquad \gamma_2=\gamma_1\left(\frac{1-T_2/T_c}{1-T_1/T_c}\right)^{n}$
\end{clave}
\begin{apoyo}\small\textbf{Ejemplo.} Etanol ($T_c=\SI{513,9}{\kelvin}$), $\gamma_1=\SI{22,0}{\milli\newton\per\meter}$ a \SI{25}{\celsius}. A \SI{50}{\celsius}:
$\displaystyle \gamma_2=22{,}0\left(\frac{1-323{,}15/513{,}9}{1-298{,}15/513{,}9}\right)^{11/9}=\SI{18,9}{\milli\newton\per\meter}$.
\end{apoyo}

\begin{alerta}\small\textbf{Errores frecuentes.} Una película de jabón tiene dos caras: $\Delta\mathcal{A}=2\ell\,\Delta x$. En Katayama--Guggenheim $T$ y $T_c$ van en kelvin, nunca en \si{\celsius}. $\gamma$ es positiva: el área tiende a disminuir sola, y aumentarla siempre requiere trabajo.\end{alerta}

%------------------------------------------------------------------ 2
\clearpage
\section{Interfases curvas: Young--Laplace}
Al minimizar el área las interfases se curvan y aparece una diferencia de presión entre el interior y el exterior de la fase. Explica gotas, aerosoles, ascenso capilar y condensación capilar.
\begin{clave}
$\displaystyle P_{\mathrm{in}}=P_{\mathrm{ext}}+\frac{2\gamma}{r}\quad\Rightarrow\quad\Delta P=\frac{2\gamma}{r}$
\end{clave}
\textbf{Gota} (líquido rodeado por su vapor) y \textbf{burbuja} (cavidad de vapor o gas en un líquido): $\Delta P=2\gamma/r$. \textbf{Pompa} (película con dos interfases): $\Delta P=4\gamma/r$.
\begin{apoyo}\small\textbf{Ejemplo.} Gota acuosa de diámetro \SI{5,0}{\micro\meter} ($\gamma=\SI{72,0}{\milli\newton\per\meter}$):
$\displaystyle\Delta P=\frac{2\,(\SI{72,0e-3}{\newton\per\meter})}{\SI{2,5e-6}{\meter}}=\SI{5,76e4}{\pascal}=\SI{57,6}{\kilo\pascal}=\SI{0,57}{atm}$.
\end{apoyo}
\begin{alerta}\small\textbf{Errores frecuentes.} Usar el diámetro en lugar del radio duplica $\Delta P$; olvidar las dos interfases de la pompa da la mitad; $\gamma$ debe ir en \si{\newton\per\meter} para obtener pascales.\end{alerta}

%------------------------------------------------------------------ 3
\section{Ecuación de Kelvin}
En partículas pequeñas la energía superficial modifica la presión de vapor ($V$ volumen molar del líquido):
\begin{clave}
$\displaystyle \ln\frac{P}{P_0}=\frac{2\gamma}{r}\,\frac{V}{RT}\qquad V=\frac{M}{\rho}$
\end{clave}
\textbf{Gota (convexo):} $P>P_0$, más cuanto menor $r$. \textbf{Poro que moja (cóncavo):} el radio cuenta como negativo y $P<P_0$ (condensación capilar). A pocos nanómetros la tensión superficial macroscópica deja de describir bien la interfase.
\begin{apoyo}\small\textbf{Ejemplos} (agua, \SI{298,15}{\kelvin}, $\gamma=\SI{71,97}{\milli\newton\per\meter}$, $V=\SI{1,807e-5}{\meter\cubed\per\mole}$).\\
\textit{Nanogota de \SI{20}{\nano\meter} de radio:}
\[\ln\frac{P}{P_0}=\frac{2\,(\SI{71,97e-3}{\newton\per\meter})(\SI{1,807e-5}{\meter\cubed\per\mole})}{(\SI{20e-9}{\meter})(\SI{8,314}{\joule\per\mole\per\kelvin})(\SI{298,15}{\kelvin})}=0{,}0525\quad\Rightarrow\quad\frac{P}{P_0}=1{,}054\]
\textit{Poro de sílice de \SI{4,0}{\nano\meter} de radio, que el agua moja} (radio negativo): $\ln\dfrac{P}{P_0}=-0{,}262\Rightarrow\dfrac{P}{P_0}=0{,}769$ ($\approx\SI{77}{\percent}$ de humedad relativa).
\end{apoyo}

%------------------------------------------------------------------ 4
\section{Capilaridad}
Intervienen tres interfases: líquido/sólido, líquido/gas y sólido/gas. El equilibrio en la línea de contacto es la ecuación de Young:
\begin{clave}
$\displaystyle\gamma_{SG}=\gamma_{SL}+\gamma_{LG}\cos\theta\quad\Rightarrow\quad\cos\theta=\frac{\gamma_{SG}-\gamma_{SL}}{\gamma_{LG}}$
\end{clave}

\begin{minipage}[t]{0.50\linewidth}\vspace{0pt}
\small
\textbf{Moja} ($0\le\theta<90^\circ$): domina la adhesión (moléculas distintas); menisco cóncavo; \emph{ascenso}.\\[2pt]
\textbf{No moja} ($90^\circ<\theta\le180^\circ$): domina la cohesión (moléculas iguales); menisco convexo; \emph{descenso} (Hg).\\[4pt]
El peso de la columna iguala a la fuerza de la tensión sobre el contorno:
$\pi r^2 g\,h\,(\rho^L-\rho^V)=2\pi r\,\gamma\cos\theta$.
\end{minipage}\hfill
\begin{minipage}[t]{0.47\linewidth}\vspace{0pt}
\centering
\begin{tikzpicture}[scale=0.78]
  % izquierda: moja
  \fill[accsoft] (0,0) rectangle (1.9,0.9);
  \draw[acc,thick] (0,0.9)--(0.45,0.9); \draw[acc,thick] (1.45,0.9)--(1.9,0.9);
  \draw[line width=1.3pt,muted] (0.45,0.2)--(0.45,2.3); \draw[line width=1.3pt,muted] (1.45,0.2)--(1.45,2.3);
  \fill[accsoft] (0.45,0.9) rectangle (1.45,1.75);
  \draw[acc,thick] (0.45,1.9) .. controls (0.8,1.55) and (1.1,1.55) .. (1.45,1.9);
  \fill[accsoft] (0.45,1.75) -- (0.45,1.9) .. controls (0.8,1.55) and (1.1,1.55) .. (1.45,1.9) -- (1.45,1.75) -- cycle;
  \draw[<->,bad,thick] (1.7,0.9)--(1.7,1.72); \node[bad,font=\scriptsize,anchor=west] at (1.72,1.3) {$h>0$};
  \node[font=\scriptsize,align=center] at (0.95,-0.35) {moja: cóncavo};
  % derecha: no moja
  \begin{scope}[xshift=3.1cm]
  \fill[accsoft] (0,0) rectangle (1.9,0.9);
  \draw[acc,thick] (0,0.9)--(0.45,0.9); \draw[acc,thick] (1.45,0.9)--(1.9,0.9);
  \draw[line width=1.3pt,muted] (0.45,0.2)--(0.45,2.3); \draw[line width=1.3pt,muted] (1.45,0.2)--(1.45,2.3);
  \fill[accsoft] (0.45,0.9) rectangle (1.45,0.45);
  \draw[acc,thick] (0.45,0.45) .. controls (0.8,0.75) and (1.1,0.75) .. (1.45,0.45);
  \fill[white] (0.46,0.46) .. controls (0.8,0.76) and (1.1,0.76) .. (1.44,0.46) -- (1.44,0.9) -- (0.46,0.9) -- cycle;
  \draw[<->,bad,thick] (1.7,0.9)--(1.7,0.5); \node[bad,font=\scriptsize,anchor=west] at (1.72,0.65) {$h<0$};
  \node[font=\scriptsize,align=center] at (0.95,-0.35) {no moja: convexo};
  \end{scope}
\end{tikzpicture}
\end{minipage}

\begin{clave}
$\displaystyle h=\frac{2\gamma\cos\theta}{g\,(\rho^L-\rho^V)\,r}\ \longrightarrow\ h=\frac{2\gamma}{r\,g\,\rho^L}\quad(\theta=0,\ \rho^L\gg\rho^V)$
\end{clave}
$h>0$ ascenso, $h<0$ descenso; $h\propto1/r$.
\begin{apoyo}\small\textbf{Ejemplo.} Capilar de microhematocrito, diámetro interno \SI{1,1}{\milli\meter}; sangre $\gamma=\SI{58}{\milli\newton\per\meter}$, $\rho=\SI{1050}{\kilo\gram\per\meter\cubed}$, $\theta=0$:
$\displaystyle h=\frac{2\,(\SI{58e-3}{\newton\per\meter})}{(\SI{0,55e-3}{\meter})(\SI{9,81}{\meter\per\second\squared})(\SI{1050}{\kilo\gram\per\meter\cubed})}=\SI{2,05e-2}{\meter}\approx\SI{2,0}{\centi\meter}$.
\end{apoyo}

%------------------------------------------------------------------ 5
\clearpage
\section{Tensión superficial en disoluciones: isoterma de Gibbs}
\begin{minipage}[t]{0.50\linewidth}\vspace{0pt}
\small
\textbf{Tipo I:} $\gamma$ sube levemente (sales inorgánicas, sacarosa); adsorción negativa.\\
\textbf{Tipo II:} $\gamma$ baja gradualmente (orgánicos solubles con grupo polar y cadena hidrocarbonada).\\
\textbf{Tipo III:} $\gamma$ cae con muy poco soluto (anfifílicos, detergentes); adsorción positiva.
\end{minipage}\hfill
\begin{minipage}[t]{0.47\linewidth}\vspace{0pt}
\centering
\begin{tikzpicture}
\begin{axis}[resumen,width=\linewidth,height=3.7cm,xlabel={$c_2$},ylabel={$\gamma$},xmin=0,xmax=1,ymin=0.2,ymax=1.15,
  xtick=\empty,ytick=\empty,axis lines=left]
  \addplot[bad,domain=0:1]{0.8+0.12*x} node[pos=0.9,above,font=\scriptsize]{I};
  \addplot[acc,domain=0:1,samples=30]{0.8-0.45*x^0.8} node[pos=0.85,above right,font=\scriptsize]{II};
  \addplot[teal,domain=0:1,samples=60]{0.8-0.5*(1-exp(-x*15))} node[pos=0.6,above,font=\scriptsize]{III};
\end{axis}
\end{tikzpicture}
\end{minipage}

El \textbf{exceso superficial} es la cantidad de soluto por unidad de área que se acumula (positivo) o falta (negativo) en la superficie, y la isoterma de Gibbs lo obtiene desde $\gamma$:
\begin{clave}
$\displaystyle \Gamma_{2(1)}=\frac{n^{\mathcal{A}}_{2(1)}}{\mathcal{A}}\quad n^{\mathcal{A}}_{2(1)}=n_2-\sum n_2^{\alpha}\qquad
\Gamma_{2(1)}=-\frac{c_2}{RT}\left(\frac{\partial\gamma}{\partial c_2}\right)_T=-\frac{1}{RT}\left(\frac{\partial\gamma}{\partial\ln c_2}\right)_T$
\end{clave}
$\partial\gamma/\partial c_2<0\Rightarrow\Gamma_{2(1)}>0$ (adsorción positiva); $\partial\gamma/\partial c_2>0\Rightarrow\Gamma_{2(1)}<0$. Área por molécula: $\sigma=1/(\Gamma_{2(1)}N_A)$.

\begin{apoyo}\small\textbf{Ejemplo 1.} Laurilsulfato de sodio a \SI{298,15}{\kelvin}, $(\partial\gamma/\partial\ln c)_T=\SI{-9,0}{\milli\newton\per\meter}$:
\[\Gamma=\frac{\SI{9,0e-3}{\newton\per\meter}}{(\SI{8,314}{\joule\per\mole\per\kelvin})(\SI{298,15}{\kelvin})}=\SI{3,63e-6}{\mole\per\meter\squared}\quad\Rightarrow\quad\sigma=\frac{1}{\Gamma N_A}=\SI{4,57e-19}{\meter\squared}=\SI{45,7}{\AA\squared}\]
\textbf{Ejemplo 2.} Polisorbato bajo su CMC: $\gamma=\gamma_0-RT\,\Gamma_{\max}\ln(1+c/a)$, $\Gamma_{\max}=\SI{2,8e-6}{\mole\per\meter\squared}$, $a=\SI{0,10}{\micro\mole\per\liter}$. Como $(\partial\gamma/\partial c)_T=-RT\Gamma_{\max}/(a+c)$, a $c=\SI{2,0}{\micro\mole\per\liter}$:
$\displaystyle\Gamma=\Gamma_{\max}\frac{c}{a+c}=\SI{2,67e-6}{\mole\per\meter\squared}$, $\sigma=\SI{62}{\AA\squared}$.
\end{apoyo}
\begin{alerta}\small\textbf{Errores frecuentes.} Con la pendiente de $\gamma$ contra $\ln c$, el factor $c$ ya está dentro de la derivada. El signo menos de la isoterma es el que hace positivo a $\Gamma$ cuando $\gamma$ disminuye.\end{alerta}

%------------------------------------------------------------------ hoja
\section{Hoja de fórmulas}
{\small\renewcommand{\arraystretch}{1.35}
\begin{tabularx}{\linewidth}{@{}>{\bfseries}l>{\raggedright\arraybackslash}Xl@{}}\toprule
Tema & Ecuación & Unidades\\\midrule
Termodinámica & $dG=\gamma\,d\mathcal{A}$;\quad $\Delta G=\gamma(\mathcal{A}_2-\mathcal{A}_1)$ & $\gamma$ en \si{\newton\per\meter}\\
Temperatura & $\gamma=\gamma_0(1-T/T_c)^{n}$,\quad $n=11/9$ & $T$ en \si{\kelvin}\\
Gota y burbuja & $\Delta P=2\gamma/r$ & \si{\pascal}\\
Pompa & $\Delta P=4\gamma/r$ & \si{\pascal}\\
Kelvin & $\ln(P/P_0)=2\gamma V/(rRT)$,\quad $V=M/\rho$ & $V$ en \si{\meter\cubed\per\mole}\\
Young & $\cos\theta=(\gamma_{SG}-\gamma_{SL})/\gamma_{LG}$ & ---\\
Altura capilar & $h=2\gamma\cos\theta\big/[g(\rho^L-\rho^V)\,r]$ & $h$ en \si{\meter}\\
Gibbs & $\Gamma_{2(1)}=-\dfrac{c_2}{RT}\left(\dfrac{\partial\gamma}{\partial c_2}\right)_T=-\dfrac{1}{RT}\left(\dfrac{\partial\gamma}{\partial\ln c_2}\right)_T$ & \si{\mole\per\meter\squared}\\
Área por molécula & $\sigma=1/(\Gamma_{2(1)}N_A)$ & \si{\meter\squared}\\\bottomrule
\end{tabularx}}

\subsection*{Guía de decisión}
\begin{multicols}{2}\small
\begin{itemize}
  \item Cambio de área de una película, trabajo: $\Delta G=\gamma\,\Delta\mathcal{A}$ (cuenta las caras).
  \item Gota, aerosol, burbuja: Laplace con $2\gamma/r$; pompa con $4\gamma/r$.
  \item Presión de vapor de una gota o de un poro: Kelvin, con el signo del radio según el menisco.
  \item Ascenso o descenso en un tubo: Young para $\theta$ y luego $h$; el signo de $\cos\theta$ manda.
  \item Soluto que baja $\gamma$ con $c$: Gibbs da $\Gamma>0$; con la pendiente en $\ln c$ no se multiplica otra vez por $c$.
  \item Antes de reemplazar: $\gamma$ en \si{\newton\per\meter}, $r$ en \si{\meter}, $T$ en \si{\kelvin}.
\end{itemize}
\end{multicols}

\end{document}
